% =============================================================================
% 3.7 Experimental Design
% =============================================================================

\section{Experimental Design}
\label{sec:experiments}

Three core experiments are defined. Each addresses one or more supporting questions. A deliberate failure injection is folded into Experiment~2 to test blast radius without requiring a separate experiment.

% ─────────────────────────────────────────────────────────────────────────────
\subsection{Experiment 1 --- VM Degradation Profile (answers SQ1)}
\label{subsec:exp1}

\paragraph{Setup.} Dagster deployed on the UTM VM with DockerRunLauncher (each job = one Docker container on the VM host). PostgreSQL 16 as the run storage backend.

\paragraph{Procedure.} Run all six workload levels (L1--L6). Each level is repeated 3 times. Record: job success rate, mean execution time, execution time standard deviation, CPU utilisation, memory utilisation, and MTTR for any failures.

\paragraph{Goal.} Produce the VM degradation curve---identifying the exact concurrency level at which success rate drops below 95\% and documenting the shape of the degradation (gradual vs.\ collapse).

% ─────────────────────────────────────────────────────────────────────────────
\subsection{Experiment 2 --- Kubernetes Isolation and Blast Radius (answers SQ2)}
\label{subsec:exp2}

\paragraph{Setup.} Kind cluster with Kubernetes Run Launcher. Identical resource profiles to Experiment~1.

\paragraph{Procedure.}

\begin{description}[style=nextline]
  \item[Part A --- Isolation comparison]
  Run identical workload levels (L1--L6) as Experiment~1, 3 repetitions each. Record same metrics.

  \item[Part B --- Blast radius test (folded in)]
  At level L4 (5 concurrent jobs), after all jobs are running, deliberately kill one container/pod mid-execution using \texttt{docker kill} (VM) or \texttt{kubectl delete pod --force} (K8s). Measure:
  \begin{enumerate}[label=(\alph*)]
    \item number of other running jobs affected,
    \item time until surviving jobs complete,
    \item success rate of surviving jobs.
  \end{enumerate}
  Repeat on VM using \texttt{docker kill} on one running job container. Compare blast radius between environments.

  \item[Part C --- Spike observation (folded in)]
  At L6 (10 concurrent jobs, exceeding single-node capacity), observe and record: how many jobs are delayed pending scheduling, how long until all jobs begin executing, any job failures due to scheduling timeout. This provides supplementary autoscaling data without a dedicated experiment.
\end{description}

\paragraph{Goal.} Calculate the delta in job success rate, execution time variance, and blast radius between VM and Kubernetes at each concurrency level.

% ─────────────────────────────────────────────────────────────────────────────
\subsection{Experiment 3 --- Overhead Measurement and Crossover Point (answers SQ3, SQ4)}
\label{subsec:exp3}

\paragraph{Setup.} Both VM and Kind environments running identical single jobs for overhead isolation, then compared across all concurrency levels.

\paragraph{Procedure.}

\begin{itemize}
  \item At each concurrency level, record pod scheduling latency and container startup time on Kubernetes (from pod event timestamps).
  \item Calculate net execution time difference between VM and Kubernetes per job per level.
  \item Plot both overhead cost (Kubernetes) and contention cost (VM) against concurrency level on the same chart.
  \item Identify the reliability crossover, performance crossover, and combined crossover point.
\end{itemize}

\paragraph{Goal.} Produce the crossover chart and identify the primary contribution---the crossover point.
